\tikzsetfigurename{12_pll-}
\chapter[Clock Generation: The Phase-Locked Loop]{Clock Generation:\\ The Phase-Locked Loop}\label{ch:pll}
\chapterlead{A crystal gives a clean 50\,MHz; the link needs a clean 2.5\,GHz in four phases. The charge-pump PLL with an LC VCO bridges the two.}

\begin{objectives}
\begin{itemize}
  \item Explain how a phase-frequency detector and charge pump turn a timing error into charge, and why the reset delay removes the dead zone.
  \item Derive the type-II loop's $\omega_n$ and $\zeta$, and design the \OS{} loop filter for a given bandwidth and phase margin.
  \item Use noise transfer functions to choose the loop bandwidth that minimizes output jitter.
  \item Size an LC tank, state the oscillation condition, and use Leeson's model to sanity-check a phase-noise number.
  \item Generate and distribute quadrature (I/Q) clocks, and budget their phase and duty-cycle errors.
\end{itemize}
\end{objectives}

%% =====================================================================
\section{The charge-pump PLL}\label{sec:cppll}

A free-running VCO drifts: its TIE random-walks (\cref{fig:ch11-tie}b) and its frequency wanders with temperature and supply. A \textbf{phase-locked loop} (PLL) fixes both by comparing the divided VCO output with a crystal reference and steering the VCO until the two are aligned. Gardner's analysis of the charge-pump PLL~\cite{p5:gardner1980} underlies this chapter; Razavi~\cite{razaviPLL}, Gardner's book~\cite{gardner2005} and ECEN~720~\cite{palermo720clk} treat it in more depth. Locked, the output frequency is exactly $N f_\mathrm{ref}$. \OS{} uses an integer-N charge-pump PLL (\cref{fig:ch12-arch}) with $f_\mathrm{ref} = 50$\,MHz and $N = 50$, reusing the team's LC-VCO PLL design~\cite{smacd2026pll} retuned from 2.4 to 2.5\,GHz.

\begin{figure}[tbp]
\centering
\begin{tikzpicture}[font=\sffamily\small,scale=0.95,transform shape]
  \node[blk,minimum width=14mm] (xo) at (0,0) {XO\\[-2pt]{\scriptsize 50\,MHz}};
  \node[blkb] (pfd) at (2.6,0) {PFD};
  \node[blkb,minimum width=13mm] (cp) at (4.9,0) {CP\\[-2pt]{\scriptsize $\pm\Icp$}};
  \node[blkb,minimum width=19mm] (vco) at (10.2,0) {LC VCO\\[-2pt]{\scriptsize $\Kvco$, 2.5\,GHz}};
  \node[blkg,minimum width=14mm] (div) at (6.4,-2.6) {$\div N$\\[-2pt]{\scriptsize $N = 50$}};
  \draw[sig] (xo) -- node[above,lbl]{$f_\mathrm{ref}$} (pfd);
  \draw[sig] (pfd.15) -- node[above,lbl]{UP} (cp.165);
  \draw[sig] (pfd.-15) -- node[below,lbl]{DN} (cp.195);
  \draw[navy,line width=0.75pt] (cp.east) -- (6.6,0) node[dot]{} -- (7.9,0) node[dot]{} -- (8.6,0);
  \draw[sig] (8.6,0) -- node[above,lbl]{$V_\mathrm{ctrl}$} (vco.west);
  \draw[color=navy] (6.6,0) to[C, l_=$C_2$] (6.6,-1.5) node[ground]{};
  \draw[color=navy] (7.9,0) to[R, l=$R$] (7.9,-0.85) to[C, l=$C_1$] (7.9,-1.7) node[ground]{};
  \node[lbl,text=bblue] at (7.25,0.45) {loop filter $Z(s)$};
  \draw[sig] (vco.east) -- (12.3,0) node[right,align=left,lbl,text=navy]{2.5\,GHz to I/Q\\generation, TX};
  \draw[sig] (11.7,0) node[dot]{} |- (div.east);
  \draw[sig] (div.west) -| node[pos=0.75,left,lbl]{$f_\mathrm{div}$} (pfd.south);
\end{tikzpicture}
\caption{\textbf{The \OS{} integer-N charge-pump PLL.} The phase-frequency detector (PFD) compares the 50\,MHz reference with the VCO divided by 50; the charge pump (CP) turns the timing difference into charge; the loop filter turns charge into the control voltage. At lock $f_\mathrm{out} = N f_\mathrm{ref} = 2.5$\,GHz.}
\label{fig:ch12-arch}
\end{figure}

\paragraph{Phase-frequency detector and charge pump}
The \textbf{PFD} is two flip-flops with their D inputs tied high (\cref{fig:ch12-pfd}). A rising reference edge sets UP, a rising divider edge sets DN, and when both are high an AND gate resets them after a short delay $T_\mathrm{rst}$. The width of UP minus the width of DN equals the timing difference $\Delta t$ between the edges. If the frequencies differ, one output dominates for many cycles, so the PFD also pulls the frequency in: it is a phase \emph{and} frequency detector. The \textbf{charge pump} switches a current $\Icp$ into the loop filter during UP and out during DN. The average current is $\Icp\,\Delta t/T_\mathrm{ref} = \Icp\,\Delta\phi/2\pi$, so the PFD-plus-CP gain is $\Icp/2\pi$ amperes per radian of reference phase.

\begin{figure}[tbp]
\centering
\begin{tikzpicture}[font=\sffamily\scriptsize,scale=0.85,transform shape,
   ff/.style={draw=navy,line width=0.7pt,fill=tile,minimum width=11mm,minimum height=11mm}]
  \node[ff] (f1) at (0,1.2) {}; \node[ff] (f2) at (0,-1.2) {};
  \foreach \f/\r in {f1/south,f2/north}{
    \node[anchor=west,text=navy] at ([yshift=2.5mm]\f.west) {D};
    \node[anchor=east,text=navy] at ([yshift=2.5mm]\f.east) {Q};
    \node[anchor=\r,text=navy,inner sep=1.5pt] at (\f.\r) {R};
    \draw[navy,line width=0.6pt] ([yshift=-4.4mm]\f.west) -- ++(0.15,0.11) -- ++(-0.15,0.11);
    \draw[navy,line width=0.7pt] ([yshift=2.5mm]\f.west) -- ++(-0.4,0) node[left]{1};}
  \draw[sig] (-1.4,0.95) node[left,text=navy]{REF} -- ([yshift=-2.5mm]f1.west);
  \draw[sig] (-1.4,-1.45) node[left,text=navy]{DIV} -- ([yshift=-2.5mm]f2.west);
  \node[and port,scale=0.7,color=navy] (and) at (2.3,0.45) {};
  \draw[navy,line width=0.7pt] (1.05,1.45) node[circ]{} |- (and.in 1);
  \draw[navy,line width=0.7pt] (1.3,-0.95) node[circ]{} |- (and.in 2);
  \node[blkg,minimum width=9mm,minimum height=5mm,font=\sffamily\scriptsize] (td) at (3.55,0.45) {$T_\mathrm{rst}$};
  \draw[navy,line width=0.7pt] (and.out) -- (td.west);
  \draw[sig] (td.east) -- (4.25,0.45) |- (0.3,0) -- (0,0) node[circ]{};
  \draw[sig] (0,0) -- (f1.south); \draw[sig] (0,0) -- (f2.north);
  \draw[sig] ([yshift=2.5mm]f1.east) -- (5.3,1.45) node[above,pos=0.85,text=navy]{UP};
  \draw[sig] ([yshift=2.5mm]f2.east) -- (5.3,-0.95) node[below,pos=0.85,text=navy]{DN};
  \draw[color=navy] (6.3,2.7) node[vcc]{} to[I, l=$\Icp$] (6.3,1.65) to[nos] (6.3,0.45) to[nos] (6.3,-0.75) to[I, l=$\Icp$] (6.3,-1.8) node[ground]{};
  \draw[dashed,gold,line width=0.6pt] (5.3,1.45) -- (6.05,1.05); \draw[dashed,gold,line width=0.6pt] (5.3,-0.95) -- (6.05,-0.15);
  \draw[sig] (6.3,0.45) node[circ]{} -- (7.5,0.45) node[right,text=navy]{to $Z(s)$};
  \node[text=navy] at (7.0,0.75) {$I_\mathrm{out}$};
\end{tikzpicture}\par\medskip
\begin{tikzpicture}[font=\sffamily\scriptsize,
   tt/.style={timing/xunit=0.245cm,timing/yunit=0.32cm,timing/slope=0.08,line width=0.75pt,draw=navy}]
  \foreach \y/\lab/\s in {0/REF/{2L8H8L8H8L8H6L},-0.62/DIV/{5L8H8L8H8L8H3L},-1.24/UP/{2L4H12L4H12L4H10L},-1.86/DN/{5L1H15L1H15L1H10L}}{
    \node[anchor=east,text=navy] at (-0.1,\y+0.16) {\lab};
    \timing[tt] at (0,\y) {\s};}
  \node[anchor=east,text=navy] at (-0.1,-2.6) {$I_\mathrm{out}$};
  \draw[navy,line width=0.75pt] (0,-2.75) -- (0.49,-2.75) -- (0.49,-2.45) -- (1.225,-2.45) -- (1.225,-2.75) -- (4.41,-2.75) -- (4.41,-2.45) -- (5.145,-2.45) -- (5.145,-2.75) -- (8.33,-2.75) -- (8.33,-2.45) -- (9.065,-2.45) -- (9.065,-2.75) -- (11.76,-2.75);
  \node[anchor=east,text=navy] at (-0.1,-3.3) {$V_\mathrm{ctrl}$};
  \draw[c2,line width=0.75pt] (0,-3.55) -- (0.49,-3.55) -- (0.49,-3.3) -- (1.225,-3.2) -- (1.225,-3.45) -- (4.41,-3.45) -- (4.41,-3.2) -- (5.145,-3.1) -- (5.145,-3.35) -- (8.33,-3.35) -- (8.33,-3.1) -- (9.065,-3.0) -- (9.065,-3.25) -- (11.76,-3.25);
  \draw[c4,<->,line width=0.6pt] (0.49,0.5) -- node[above]{$\Delta t$} (1.225,0.5);
  \draw[muted,dotted] (0.49,0.45) -- (0.49,-3.6); \draw[muted,dotted] (1.225,0.45) -- (1.225,-3.6);
  \draw[golddk,->,line width=0.6pt] (2.9,-2.25) node[right]{both on for $T_\mathrm{rst}$: currents cancel, no dead zone} -- (1.45,-1.75);
  \node[lbl,anchor=west] at (5.4,-3.65) {jumps by $\Icp R$ while pumping; net step $\Icp\Delta t/C_1$};
\end{tikzpicture}
\caption{\textbf{The PFD measures the edge offset as a pulse width; the charge pump turns it into charge.} Top: two reset flip-flops, an AND gate with reset delay $T_\mathrm{rst}$, and the up/down current sources. Bottom: REF leads DIV by $\Delta t$, so UP is wider than DN by $\Delta t$; both are on together for $T_\mathrm{rst}$ and cancel. Each cycle the control voltage jumps by $\Icp R$ while current flows and keeps a net step of $\Icp\Delta t/C_1$.}
\label{fig:ch12-pfd}
\end{figure}

\paragraph{Dead zone and mismatch}
Without the reset delay, a tiny $\Delta t$ would give UP and DN pulses too short to switch the charge pump on, so near lock the loop would get no correction: a \textbf{dead zone} in which the VCO drifts. With $T_\mathrm{rst}$ of 100--200\,ps both pulses always last at least $T_\mathrm{rst}$, and the net charge stays proportional to $\Delta t$ down to zero. The price is mismatch: if the currents differ by $\Delta I$, the loop locks at a static offset $\Delta t = (\Delta I/\Icp)\,T_\mathrm{rst}$ (7.5\,ps for 5\% and 150\,ps), and the current pulses every 20\,ns create a \textbf{reference spur}, a periodic jitter that $C_2$ and the loop must filter.

%% =====================================================================
\section{Loop dynamics}\label{sec:pll-loop}

Near lock every block is linear in \emph{phase}, which gives the model of \cref{fig:ch12-linmodel}. The PFD and CP give $\Icp/2\pi$ (A/rad); the loop filter impedance $Z(s)$ turns current into $V_\mathrm{ctrl}$; the VCO integrates frequency into phase, $\Kvco/s$ (rad/s/V divided by $s$); and the divider scales phase by $1/N$. The same figure marks where each noise source enters.

\begin{figure}[tbp]
\centering
\begin{tikzpicture}[font=\sffamily\small,
  b/.style={blkb,minimum width=13mm,minimum height=8mm},
  sm/.style={sumnode,minimum size=5mm,font=\scriptsize},
  nz/.style={font=\sffamily\scriptsize,text=c4,align=center},
  tshape/.style={line width=0.7pt,draw=golddk}]
  \node[text=navy] (in) at (0,0) {$\phi_\mathrm{ref}$};
  \node[sm] (s0) at (1.15,0) {$+$};
  \node[sm] (s1) at (2.3,0) {$\Sigma$};
  \node[b] (kp) at (3.75,0) {$\dfrac{\Icp}{2\pi}$};
  \node[sm] (s2) at (5.15,0) {$+$};
  \node[b] (zf) at (6.45,0) {$Z(s)$};
  \node[sm] (s3) at (7.75,0) {$+$};
  \node[b] (kv) at (9.1,0) {$\dfrac{\Kvco}{s}$};
  \node[sm] (s4) at (10.5,0) {$+$};
  \node[text=navy] (out) at (12.0,0) {$\phi_\mathrm{out}$};
  \node[blkg,minimum width=11mm,minimum height=7mm] (dv) at (6.45,-1.35) {$1/N$};
  \draw[sig] (in) -- (s0); \draw[sig] (s0) -- (s1); \draw[sig] (s1) -- (kp); \draw[sig] (kp) -- (s2);
  \draw[sig] (s2) -- (zf); \draw[sig] (zf) -- (s3); \draw[sig] (s3) -- (kv); \draw[sig] (kv) -- (s4); \draw[sig] (s4) -- (out);
  \draw[sig] (11.25,0) node[dot]{} |- (dv.east); \draw[sig] (dv.west) -| node[pos=0.95,left,font=\scriptsize]{$-$} (s1.south);
  \foreach \s/\t/\k in {s0/{$\phi_{n,\mathrm{ref}}$}/LP,s2/{$i_{n,\mathrm{CP}}$}/LP,s3/{$v_{n,R}$}/BP,s4/{$\phi_{n,\mathrm{VCO}}$}/HP}{
    \node[nz] (n\s) at ([yshift=12mm]\s) {\t};
    \draw[sig,draw=c4] (n\s) -- (\s);}
  % tiny transfer-shape icons
  \begin{scope}[shift={(1.55,1.05)},x=0.45cm,y=0.3cm]
    \draw[tshape] (0,0.8) -- (0.6,0.8) -- (1,0); \node[font=\sffamily\tiny,text=golddk] at (0.5,-0.4) {LP $\times N$};\end{scope}
  \begin{scope}[shift={(5.55,1.05)},x=0.45cm,y=0.3cm]
    \draw[tshape] (0,0.8) -- (0.6,0.8) -- (1,0); \node[font=\sffamily\tiny,text=golddk] at (0.5,-0.4) {LP};\end{scope}
  \begin{scope}[shift={(8.15,1.05)},x=0.45cm,y=0.3cm]
    \draw[tshape] (0,0) -- (0.35,0.8) -- (0.65,0.8) -- (1,0); \node[font=\sffamily\tiny,text=golddk] at (0.5,-0.4) {BP};\end{scope}
  \begin{scope}[shift={(10.9,1.05)},x=0.45cm,y=0.3cm]
    \draw[tshape] (0,0) -- (0.4,0.8) -- (1,0.8); \node[font=\sffamily\tiny,text=golddk] at (0.5,-0.4) {HP};\end{scope}
\end{tikzpicture}
\caption{\textbf{Linear phase model of the CP-PLL, with its noise sources.} Each source reaches the output through a different filter: reference and charge-pump noise are low-pass filtered (and multiplied by $N$), loop-filter resistor noise is band-pass, and VCO noise is high-pass. The loop bandwidth sets the crossover between ``trust the reference'' and ``trust the VCO''.}
\label{fig:ch12-linmodel}
\end{figure}

\paragraph{Type-II second-order loop} Ignore $C_2$ for now, so $Z(s) = R + 1/(sC_1)$. The open-loop gain is
\[
  LG(s) = \frac{\Icp}{2\pi}\,Z(s)\,\frac{\Kvco}{s}\,\frac1N = \frac{\Icp\Kvco}{2\pi N C_1}\cdot\frac{1 + sRC_1}{s^2} .
\]
It has two integrators (the VCO and $C_1$), so it is a \textbf{type-II} loop: it settles to zero phase error even after a frequency step. The closed-loop transfer $\phi_\mathrm{out}/\phi_\mathrm{ref} = N\,LG/(1+LG)$ becomes, after multiplying through by $s^2$,
\begin{keyeq}[CP-PLL natural frequency and damping]
\begin{equation}
  \frac{\phi_\mathrm{out}}{\phi_\mathrm{ref}} = N\,\frac{2\zeta\omega_n s + \omega_n^2}{s^2 + 2\zeta\omega_n s + \omega_n^2},
  \qquad
  \omega_n = \sqrt{\frac{\Icp\Kvco}{2\pi N C_1}}, \qquad
  \zeta = \frac{R}{2}\sqrt{\frac{\Icp\Kvco C_1}{2\pi N}} .
  \label{eq:ch12-wnz}
\end{equation}
\end{keyeq}
Matching coefficients gives $\omega_n^2 = \Icp\Kvco/(2\pi N C_1)$ and $2\zeta\omega_n = \Icp\Kvco R/(2\pi N)$. Raising $R$ raises the damping without changing $\omega_n$; raising $C_1$ lowers $\omega_n$ and raises $\zeta$. The zero in the numerator causes some \textbf{jitter peaking} at every damping (\cref{fig:ch12-resp}, \cref{tab:ch12-zeta}): the loop amplifies reference jitter near $\omega_n$. The $-3$\,dB bandwidth is $\omega_{-3\mathrm{dB}} = \omega_n\big[1 + 2\zeta^2 + \sqrt{(1+2\zeta^2)^2 + 1}\,\big]^{1/2}$, which tends to $2\zeta\omega_n$ for large $\zeta$.

\begin{table}[htbp]
\centering
\caption{Bandwidth and peaking of the type-II loop of \cref{eq:ch12-wnz} (computed).}
\label{tab:ch12-zeta}
\small
\begin{tabular}{lcccccc}
\toprule
$\zeta$ & 0.5 & 0.707 & 1 & 2 & 4 & 4.5 \\
\midrule
$\omega_{-3\mathrm{dB}}/\omega_n$ & 1.82 & 2.06 & 2.48 & 4.25 & 8.13 & 9.11 \\
Peaking (dB) & 3.3 & 2.1 & 1.25 & 0.40 & 0.12 & 0.09 \\
\bottomrule
\end{tabular}
\end{table}

\begin{figure}[tbp]
\centering
\includegraphics[width=\linewidth]{ch12_resp}
\caption{\textbf{Damping trades peaking and overshoot for bandwidth.} (a) Closed-loop $|H|/N$: low damping gives a sharp, peaky response; high damping a wide, flat one. (b) Output phase after a reference phase step: $\zeta = 0.5$ overshoots by 30\%, $\zeta = 1$ by 13\%. Transmit PLLs use $\zeta \approx 0.7$--1; CDRs, which must not amplify jitter, use $\zeta \gg 1$ (\cref{ch:cdr}).}
\label{fig:ch12-resp}
\end{figure}

\paragraph{The third pole and phase margin} The capacitor $C_2$ in parallel smooths the voltage step $\Icp R$ seen in \cref{fig:ch12-pfd}, which would otherwise modulate the VCO every reference cycle. It adds a pole at $\omega_{p3} = 1/(R C_s)$ with $C_s = C_1C_2/(C_1+C_2)$, and the loop becomes third order. The usual design places the unity-gain frequency $\omega_u$ at the geometric mean of the zero $\omega_z = 1/(RC_1)$ and $\omega_{p3}$, where the phase margin $\Phi_m$ is largest. With $K_C = C_1/C_2$:
\begin{equation}
\begin{gathered}
  K_C = 2\tan^2\Phi_m + 2\tan\Phi_m\sqrt{\tan^2\Phi_m + 1}, \qquad
  \omega_z = \frac{\omega_u}{\sqrt{1+K_C}}, \qquad
  \omega_{p3} = \omega_u\sqrt{1+K_C}, \\
  C_2 = \frac{\Icp\Kvco}{2\pi N\omega_u^2}\sqrt{\frac{\omega_u^2+\omega_z^2}{\omega_u^2+\omega_{p3}^2}}, \qquad C_1 = K_C C_2, \qquad R = \frac{1}{\omega_z C_1} .
\end{gathered}
  \label{eq:ch12-design}
\end{equation}
For $\Phi_m = 60^\circ$, $K_C = 12.9$. Keep $f_u \le f_\mathrm{ref}/10$, or the sampled nature of the PFD makes the continuous-time model wrong.

\begin{example}{Loop filter for the Open-SERDES PLL}\label{we:ch12-lf}
\textbf{Given:} $f_\mathrm{ref} = 50$\,MHz, $N = 50$, $\Kvco = 2\pi\times150$\,MHz/V, $\Icp = 200$\,\textmu{}A, $f_u = 1$\,MHz, $\Phi_m = 60^\circ$.

\textbf{Pole and zero:} $K_C = 12.93$, so $f_z = 1\,\mathrm{MHz}/\sqrt{13.93} = 268$\,kHz and $f_{p3} = 3.73$\,MHz.
\textbf{Components} from \cref{eq:ch12-design}: $C_2 = \mathbf{4.07}$\,pF, $C_1 = K_C C_2 = \mathbf{52.6}$\,pF, $R = 1/(\omega_z C_1) = \mathbf{11.3}$\,k\Ohm.
\textbf{Check} with \cref{eq:ch12-wnz}: $f_n = 537$\,kHz and $\zeta = 1.00$. The full third-order loop, simulated in Python (\cref{fig:ch12-loop}), crosses unity gain at 1.00\,MHz with 60.0$^\circ$ margin and has $f_{-3\mathrm{dB}} = 1.56$\,MHz with 1.7\,dB peaking.
\textbf{Practical:} $f_\mathrm{ref}/f_u = 50$, comfortably above 10. $C_1$ is 0.035\,mm$^2$ of MIM at about 1.5\,fF/\textmu{}m$^2$ (check the PDK value), the largest item in the PLL layout. $R$ adds $\sqrt{4kTR} = 13.7$\,nV/$\sqrt{\mathrm{Hz}}$ of band-pass noise.
\end{example}

\begin{figure}[tbp]
\centering
\includegraphics[width=\linewidth]{ch12_loop}
\caption{\textbf{The designed loop has 60$^\circ$ of phase margin at 1\,MHz.} (a) Open-loop gain and phase of \cref{we:ch12-lf}: the phase is closest to $-180^\circ$ furthest from it at $f_u$, midway between $f_z$ and $f_{p3}$ on a log scale. (b) The third-order closed loop (navy) against the second-order fit with $\zeta = 1$: they agree up to the bandwidth, and $C_2$ adds extra roll-off above it.}
\label{fig:ch12-loop}
\end{figure}

%% =====================================================================
\section{PLL noise and the optimum bandwidth}\label{sec:ch12-noise}

From \cref{fig:ch12-linmodel}, each noise source sees its own transfer function:
\begin{equation}
  \mathrm{NTF}_\mathrm{ref} = N\frac{LG}{1+LG}, \quad
  \mathrm{NTF}_\mathrm{CP} = \frac{2\pi}{\Icp}\,N\frac{LG}{1+LG}, \quad
  \mathrm{NTF}_\mathrm{VCO} = \frac{1}{1+LG}, \quad
  \mathrm{NTF}_R = \frac{\Kvco/s}{1+LG} .
\end{equation}
Reference and charge-pump noise pass at low frequency and are multiplied by $N$ ($+34$\,dB for $N = 50$); VCO noise is suppressed at low frequency and passes above the bandwidth (\cref{fig:ch12-noise}a). Widening the loop therefore lets in more reference noise and keeps out more VCO noise. The best bandwidth is roughly where the free-running VCO noise crosses the in-band floor.

\begin{figure}[tbp]
\centering
\includegraphics[width=\linewidth]{ch12_noise}
\caption{\textbf{The loop bandwidth trades reference noise against VCO noise.} (a) Noise transfer for $f_u$ = 1 and 3\,MHz: the reference and CP path is low-pass, the VCO path high-pass. (b) Integrated output jitter (10\,kHz to 1.25\,GHz) against $f_u$, for the VCO of \cref{we:ch11-pn} and a $-110$\,dBc/Hz in-band floor. The total is flat near its 3\,MHz minimum of 1.01\,ps. If the VCO is 22\,dB quieter (\cref{sec:lcvco}), the optimum moves below 0.3\,MHz (dotted).}
\label{fig:ch12-noise}
\end{figure}

\begin{example}{Choosing the PLL bandwidth}\label{we:ch12-bw}
\textbf{Given:} the VCO noise of \cref{we:ch11-pn} ($-100.8$\,dBc/Hz at 1\,MHz), and an in-band floor of $-110$\,dBc/Hz at the output. That floor corresponds to a normalized PLL figure of merit of $-110 - 10\log_{10}(50\,\mathrm{MHz}) - 20\log_{10}50 = -221$\,dBc/Hz, typical of a good CP-PLL.

\textbf{Sweep} (60$^\circ$ loops of \cref{eq:ch12-design}, \cref{fig:ch12-noise}b): $f_u$ = 0.5, 1, 2, 3, 6\,MHz gives 1.91, 1.34, 1.06, \textbf{1.01}, 1.14\,ps rms. At 3\,MHz the VCO and in-band terms are equal (0.71 and 0.72\,ps).

\textbf{Redesign at 3\,MHz:} capacitors scale as $\Icp/f_u^2$ and $R$ as $f_u/\Icp$. Tripling $\Icp$ to 600\,\textmu{}A keeps $R = 11.3$\,k\Ohm{} and gives $C_1 = 17.5$\,pF, $C_2 = 1.36$\,pF (a third of the area), $f_{-3\mathrm{dB}} = 4.7$\,MHz and $f_\mathrm{ref}/f_u = 17$. The resistor noise adds only 0.11\,ps in RSS. (Keeping 200\,\textmu{}A would need $R = 34$\,k\Ohm{} and a 0.45\,pF $C_2$ that parasitics would swamp.)
\end{example}

\begin{insight}
Choose $f_u \approx 3$\,MHz with $\Icp = 600$\,\textmu{}A. The minimum is shallow (within 15\% from 2 to 6\,MHz), so bandwidth can be traded for area or settling time. But it depends on the VCO's free-running noise, which is uncertain (\cref{sec:lcvco}): pin that down before freezing $f_u$.
\end{insight}

%% =====================================================================
\section{The LC VCO}\label{sec:lcvco}

An LC tank stores energy and loses a little of it each cycle in the tank resistance. The \textbf{cross-coupled pair} of \cref{fig:ch12-vco} replaces that loss: looking into the two drains, the pair presents a negative resistance $-2/\gm$. The tank's loss is a parallel resistance $R_p = Q\omega_0 L$ (differential), so
\begin{keyeq}[LC oscillator: frequency and start-up]
\begin{equation}
  f_0 = \frac{1}{2\pi\sqrt{L\,C_\mathrm{tot}}}, \qquad
  R_p = Q\,\omega_0 L, \qquad
  \text{oscillation if}\ \ \frac{\gm R_p}{2} \ge 1 \ \ (\text{design for 2--3}).
  \label{eq:ch12-lc}
\end{equation}
\end{keyeq}
The amplitude grows until the pair limits, at about $(4/\pi)\,I_\mathrm{tail}R_p/2$ (differential). $C_\mathrm{tot}$ sums the varactors (fine tuning, which sets $\Kvco$), a switched-capacitor bank (coarse bands for process corners) and all parasitics.

\begin{figure}[tbp]
\centering
\begin{tikzpicture}[font=\sffamily\footnotesize,scale=0.9,transform shape]
  \draw[color=navy] (0,4.3) to[L, l=$L/2$] (2.2,4.3) to[L, l=$L/2$] (4.4,4.3);
  \draw[color=navy] (2.2,4.3) -- (2.2,4.9) node[vcc]{$V_\mathrm{DD}$};
  \draw[color=navy] (0,4.3) -- (0,0.9); \draw[color=navy] (4.4,4.3) -- (4.4,0.9);
  \draw[color=navy] (0,3.0) node[circ]{} to[vC, l=$C_v$] (2.2,3.0) to[vC, l=$C_v$] (4.4,3.0) node[circ]{};
  \draw[color=navy] (2.2,3.0) -- (2.2,2.55) node[right,text=navy]{$V_\mathrm{ctrl}$};
  \draw[color=navy] (0,1.85) node[circ]{} to[C, l=$C_s$] (1.4,1.85) to[nos, l_=$b_k$] (3.0,1.85) to[C, l=$C_s$] (4.4,1.85) node[circ]{};
  \node[nmos,xscale=-1,anchor=D] (m1) at (0,0.9) {}; \node[nmos,anchor=D] (m2) at (4.4,0.9) {};
  \node[text=navy,anchor=east] at ([xshift=-1mm]m1.center) {M1}; \node[text=navy,anchor=west] at ([xshift=1mm]m2.center) {M2};
  \draw[color=navy] (m1.G) -- ++(0.45,0) -- (3.0,1.3) -- (4.4,1.3) node[circ]{};
  \draw[color=navy] (m2.G) -- ++(-0.45,0) -- (1.4,1.3) -- (0,1.3) node[circ]{};
  \draw[color=navy] (m1.S) |- (2.2,-1.0) -| (m2.S);
  \draw[color=navy] (2.2,-1.0) node[circ]{} to[I, l=$I_\mathrm{tail}$] (2.2,-2.4) node[ground]{};
  \draw[color=navy] (0,3.65) node[circ]{} -- (-0.6,3.65) node[left,text=navy]{$V_\mathrm{out}^{+}$};
  \draw[color=navy] (4.4,3.65) node[circ]{} -- (5.0,3.65) node[right,text=navy]{$V_\mathrm{out}^{-}$};
  \draw[region] (-0.3,1.55) rectangle (4.7,4.65);
  \node[lbl,anchor=west,align=left] at (5.0,2.4) {tank: $L$ = 4\,nH, $Q$ = 16.8,\\loss $R_p = Q\omega_0L$ = 1.06\,k\Ohm};
  \node[lbl,anchor=west,align=left] at (5.0,0.0) {cross-coupled pair:\\$R_\mathrm{in} = -2/\gm$};
\end{tikzpicture}
\caption{\textbf{Cross-coupled LC VCO.} A center-tapped inductor and the capacitors set $f_0$; varactors $C_v$ give fine tuning through $V_\mathrm{ctrl}$, and switched capacitors $C_s$ (bits $b_k$) select coarse bands. The cross-coupled pair cancels the tank loss when $\gm R_p/2 \ge 1$, and the tail current sets the amplitude. Drawn with MOSFETs; npn13G2 HBTs work the same way, with capacitive cross-coupling to keep $V_\mathrm{CE}$ below BV$_\mathrm{CEO}$.}
\label{fig:ch12-vco}
\end{figure}

\begin{example}{The 2.5\,GHz tank}\label{we:ch12-tank}
\textbf{Given:} the team's 4\,nH differential inductor with $Q = 16.8$~\cite{smacd2026pll}, $f_0 = 2.5$\,GHz.
\textbf{Capacitance:} $C_\mathrm{tot} = 1/(\omega_0^2L) = 1/\big((2\pi\cdot2.5\times10^9)^2\cdot4\times10^{-9}\big) = \mathbf{1.01}$\,pF, including all parasitics.
\textbf{Loss:} $R_p = 16.8\times2\pi\cdot2.5\times10^9\times4\times10^{-9} = \mathbf{1.06}$\,k\Ohm.
\textbf{Start-up:} $\gm \ge 2/R_p = 1.89$\,mS; with a 3$\times$ margin, 5.7\,mS per device (about 0.15\,mA of collector current for an HBT, more for a MOSFET).
\textbf{Tuning:} $\pm10\%$ in frequency needs $C_\mathrm{max}/C_\mathrm{min} = (1.1/0.9)^2 = 1.49$ across varactors and bank together, with the fixed parasitics counted in both.
\end{example}

\paragraph{Phase noise: Leeson's model}
Leeson's empirical model~\cite{p5:leeson1966} captures the three regions of an oscillator's phase noise (\cref{fig:ch12-leeson}a):
\begin{equation}
  \mathcal{L}(\Delta f) = 10\log_{10}\!\left[\frac{2FkT}{P_\mathrm{sig}}\left(1 + \Big(\frac{f_0}{2Q\,\Delta f}\Big)^2\right)\left(1 + \frac{f_{1/f^3}}{\Delta f}\right)\right] ,
  \label{eq:ch12-leeson}
\end{equation}
where $P_\mathrm{sig} = V_\mathrm{pk}^2/(2R_p)$ is the power in the tank and $F$ an excess-noise factor of the active devices. Thermal noise near the carrier is shaped by the tank into a $1/\Delta f^2$ slope, flicker noise adds $1/\Delta f^3$ close in, and far out the noise is flat. For a fixed inductor, $\mathcal{L} \propto 1/(Q\,V_\mathrm{pk}^2)$: a higher $Q$ and a larger swing both help (\cref{fig:ch12-leeson}b).

\begin{figure}[tbp]
\centering
\includegraphics[width=\linewidth]{ch12_leeson}
\caption{\textbf{Leeson's model says the team's tank should be about 21\,dB quieter than the quoted number.} (a) Leeson's model for the 2.5\,GHz tank ($Q = 16.8$, 0.6\,V peak, $F = 2$): $-122$\,dBc/Hz at 1\,MHz. Matching $-100.8$\,dBc/Hz would need 22\,dB of excess noise (dashed). (b) $\mathcal{L}$(1\,MHz) against tank amplitude: even 50\,mV of swing would beat the quoted value.}
\label{fig:ch12-leeson}
\end{figure}

\begin{pitfall}[Is $-100.8$\,dBc/Hz the VCO, or the PLL?]
With $Q = 16.8$, $R_p = 1.06$\,k\Ohm, 0.6\,V of swing and $F = 2$, \cref{eq:ch12-leeson} predicts $-122$\,dBc/Hz at 1\,MHz. The common oscillator figure of merit,
\[\mathrm{FoM} = -\mathcal{L}(\Delta f) + 20\log_{10}(f_0/\Delta f) - 10\log_{10}(P/1\,\mathrm{mW}),\]
would be only 162\,dB for $-100.8$\,dBc/Hz at 5\,mW of VCO power, against 180--190\,dB for typical LC VCOs. The quoted value is probably the \emph{closed-loop} PLL output near the loop bandwidth, or the VCO has excess noise (tail noise, low swing). This book uses it as free-running VCO noise, the pessimistic reading; resolving it is a Phase~0 task, because it moves the optimum of \cref{we:ch12-bw}.
\end{pitfall}

\begin{insight}[Retuning the team's 2.4\,GHz design to 2.5\,GHz]
Keep the 4\,nH inductor: its $Q$ and layout are proven. The tank capacitance must fall from 1.10 to 1.01\,pF, by 86\,fF or 8\%: remove one switched-capacitor unit or shift the band code, and re-check that $V_\mathrm{ctrl}$ sits mid-range at all corners. $R_p$ rises by 4\% with frequency, so the start-up margin improves slightly. Replace the fractional-N divider with an integer $\div50$ and redesign the loop filter for $f_u \approx 3$\,MHz (\cref{we:ch12-bw}). Target $\Kvco \approx 2\pi\times150$\,MHz/V, so that supply and control noise stay small, and cover PVT with a 3--5 bit capacitor bank.
\end{insight}

%% =====================================================================
\section{Quadrature clocks and distribution}\label{sec:quadrature}

The half-rate receiver samples data on the 0$^\circ$ and 180$^\circ$ edges of the 2.5\,GHz clock and the data transitions on the 90$^\circ$ and 270$^\circ$ edges, with four samplers D0, E0, D1, E1 (\cref{sec:sampler}). It needs four phases, I, Q, $\bar{\mathrm{I}}$ and $\bar{\mathrm{Q}}$, spaced by exactly 100\,ps. \OS{} derives them from the 2.5\,GHz VCO with an \textbf{RC polyphase filter} (PPF); the alternative is a 5\,GHz VCO followed by a $\div2$ (\cref{fig:ch12-quad}).

\begin{figure}[tbp]
\centering
\begin{subfigure}[b]{0.62\linewidth}\centering
\begin{tikzpicture}[font=\sffamily\scriptsize,scale=0.72,transform shape]
  \foreach \i/\y in {0/0,1/-1.25,2/-2.5,3/-3.75}{
    \draw[color=navy] (0.6,\y) to[R] (3.0,\y) -- (3.8,\y) to[R] (6.2,\y) -- (6.8,\y);}
  \foreach \y/\yn in {0/-1.25,-1.25/-2.5,-2.5/-3.75}{
    \draw[color=navy] (1.0,\y) node[circ]{} to[C] (2.6,\yn) node[circ]{};
    \draw[color=navy] (4.2,\y) node[circ]{} to[C] (5.8,\yn) node[circ]{};}
  \draw[color=navy] (1.0,-3.75) node[circ]{} to[C] (2.6,-5.0) -- (3.0,-5.0) node[right,text=muted]{to 0$^\circ$};
  \draw[color=navy] (4.2,-3.75) node[circ]{} to[C] (5.8,-5.0) -- (6.2,-5.0) node[right,text=muted]{to 0$^\circ$};
  \foreach \t/\y in {{$V_\mathrm{in}^{+}$}/0,{$V_\mathrm{CM}$}/-1.25,{$V_\mathrm{in}^{-}$}/-2.5,{$V_\mathrm{CM}$}/-3.75}{\node[left,text=navy] at (0.6,\y) {\t};}
  \foreach \t/\y in {I/0,Q/-1.25,{$\bar{\mathrm{I}}$}/-2.5,{$\bar{\mathrm{Q}}$}/-3.75}{\node[right,text=navy] at (6.8,\y) {\t};}
  \node[lbl] at (1.8,0.7) {stage 1: $R_1, C_1$, 2.0\,GHz}; \node[lbl] at (5.0,0.7) {stage 2: $R_2, C_2$, 3.1\,GHz};
\end{tikzpicture}
\caption{Two-stage RC polyphase filter}
\end{subfigure}\hfill
\begin{subfigure}[b]{0.36\linewidth}\centering
\begin{tikzpicture}[font=\sffamily\scriptsize]
  \node[blkb,minimum width=12mm,minimum height=10mm,font=\sffamily\scriptsize] (l1) at (0,0) {D-latch\\CK};
  \node[blkb,minimum width=12mm,minimum height=10mm,font=\sffamily\scriptsize] (l2) at (2.2,0) {D-latch\\$\overline{\mathrm{CK}}$};
  \draw[sig] (l1.east) -- (l2.west);
  \draw[sig] (l2.north) -- ++(0,0.45) -| node[pos=0.25,above]{$\bar{\mathrm{Q}}$ (inverted)} ([xshift=-4mm]l1.west) -- (l1.west);
  \draw[sig] (1.1,0) node[dot]{} -- (1.1,-0.85) -- (3.4,-0.85) node[right]{I};
  \draw[sig] (l2.east) -- (3.4,0) node[right]{Q};
  \draw[clk] (-0.9,-1.45) node[left,text=golddk]{5\,GHz} -- (2.2,-1.45);
  \draw[clk] (0,-1.45) -- (l1.south); \draw[clk] (2.2,-1.45) -- (l2.south);
\end{tikzpicture}
\caption{Alternative: $\div2$ from a 5\,GHz VCO}
\end{subfigure}
\caption{\textbf{Two ways to make quadrature.} (a) Each PPF stage has an $R$ from each input phase and a $C$ from the previous phase (the last capacitor wraps around to the 0$^\circ$ output). Driven differentially, one stage gives outputs exactly 90$^\circ$ apart, with equal amplitude only at $f = 1/(2\pi RC)$; a second stage with a staggered pole keeps them equal across process. (b) A master-slave $\div2$ gives I and Q at 2.5\,GHz from both edges of a 5\,GHz clock; its accuracy depends on the 5\,GHz duty cycle, and it needs a VCO at twice the frequency.}
\label{fig:ch12-quad}
\end{figure}

\paragraph{How accurate is the polyphase filter?}
For a single stage driven differentially, the outputs are $V/(1 + j\omega RC)$ and $j\omega RC\,V/(1 + j\omega RC)$: always 90$^\circ$ apart, with amplitude ratio $\omega RC$. A $\pm20\%$ process shift of $RC$ therefore leaves 1.6--1.9\,dB of amplitude imbalance, while two stages with poles at 2.0 and 3.1\,GHz keep it below 0.07\,dB (\cref{fig:ch12-poly}a). Phase errors come from random mismatch instead: with 1\% mismatch on every $R$ and $C$, Monte Carlo gives 0.57$^\circ$ rms for one stage and 0.59$^\circ$ for two (\cref{fig:ch12-poly}b). The last stage sets the phase error, so it needs large, well-matched resistors and MIM capacitors in a common-centroid layout. The PPF is passive and lossy, $-5.2$\,dB for two unloaded stages, so CML buffers must restore the swing.

\begin{figure}[tbp]
\centering
\includegraphics[width=\linewidth]{ch12_poly}
\caption{\textbf{Two PPF stages fix the amplitude; matching sets the phase.} (a) I/Q amplitude imbalance against frequency. One stage is balanced only at its pole; two staggered stages stay within 0.07\,dB over the whole $\pm20\%$ $RC$ range (shaded) at 2.5\,GHz. (b) Quadrature phase error over 3000 Monte Carlo runs with 1\% $R$ and $C$ mismatch: about 0.6$^\circ$ rms (0.66\,ps) either way.}
\label{fig:ch12-poly}
\end{figure}

\paragraph{What the errors do to the half-rate samplers}
At 2.5\,GHz one degree is 1.11\,ps. The CDR locks the edge samplers (Q, $\bar{\mathrm{Q}}$) onto the data transitions, so an I/Q error $\delta$ moves the data samplers $\delta\times1.11$\,ps away from the eye centre: a static offset that eats timing margin. A duty-cycle error is worse, because the D1 and E1 samplers fire on the falling edges: a 1\% duty error shifts them by 4\,ps relative to D0 and E0, a two-valued DCD seen by every other bit.

\paragraph{Distribution} CML limiting buffers after the PPF restore the swing, remove the residual imbalance and keep the slew rate high (\cref{eq:ch11-slope}). Give them 2--3$\times f_\mathrm{clk}$ of bandwidth (shunt peaking or tuned loads): a buffer near $f_\mathrm{clk}$ amplifies jitter and duty-cycle error. A duty-cycle correction (DCC) loop that servoes the clock's average to zero holds the duty cycle within about $\pm0.5\%$.

\begin{example}{The RX clock-path line of the jitter budget}\label{we:ch12-iq}
\textbf{I/Q error:} two-stage PPF, 1\% mismatch: $3\sigma = 3\times0.59^\circ = 1.8^\circ \to$ \textbf{2.0\,ps}.
\textbf{Duty cycle:} $\pm0.5\%$ after DCC: $0.5\% \times 400$\,ps $\to$ \textbf{2.0\,ps}.
\textbf{PI nonlinearity:} $\pm0.6$\,ps with corrected weights (\cref{sec:picdr}) $\to$ \textbf{1.2\,ps}.
\textbf{Total} $5.2$\,ps, inside the 6\,ps allocated in \cref{tab:ch11-budget}, with 0.8\,ps of margin. Without DCC, a 52\% duty cycle alone would cost 8\,ps.
\end{example}

\begin{specnote}[Open-SERDES clock generation targets]
Integer-N CP-PLL, $f_\mathrm{ref} = 50$\,MHz, $N = 50$, $\Kvco \approx 2\pi\times150$\,MHz/V, $\Icp = 600$\,\textmu{}A, $f_u \approx 3$\,MHz ($f_{-3\mathrm{dB}} \approx 4.7$\,MHz), 60$^\circ$ phase margin. Jitter $\le 1.0$\,ps rms from 10\,kHz and $\le 0.6$\,ps above 4\,MHz. Quadrature from a two-stage PPF and CML buffers: I/Q error $\le 2^\circ$ ($3\sigma$), duty cycle $50 \pm 0.5\%$. Power estimate: PLL about 13\,mW (the team's design draws 12.73\,mW); PPF buffers and distribution 10--15\,mW (four CML stages of about 1.5\,mA at 2.5\,V).
\end{specnote}

\begin{takeaways}
\begin{itemize}
  \item The PFD turns the edge offset into an UP/DN width difference; the charge pump gives $\Icp/2\pi$ A/rad. The reset delay removes the dead zone; current mismatch causes a static offset and reference spurs.
  \item Type-II loop: $\omega_n = \sqrt{\Icp\Kvco/(2\pi NC_1)}$, $\zeta = (R/2)\sqrt{\Icp\Kvco C_1/(2\pi N)}$; peaking is 2.1\,dB at $\zeta = 0.707$ and 0.4\,dB at $\zeta = 2$.
  \item For 60$^\circ$ margin at 1\,MHz: $C_2$ = 4.07\,pF, $C_1$ = 52.6\,pF, $R$ = 11.3\,k\Ohm{} ($f_{-3\mathrm{dB}}$ = 1.56\,MHz); at the 3\,MHz optimum, $\Icp$ = 600\,\textmu{}A, $C_1$ = 17.5\,pF.
  \item Reference noise is low-pass and VCO noise high-pass; the bandwidth that equalizes them, about 3\,MHz here, minimizes jitter at 1.0\,ps rms.
  \item The 4\,nH, $Q = 16.8$ tank needs 1.01\,pF at 2.5\,GHz, has $R_p$ = 1.06\,k\Ohm{}, and oscillates if $\gm R_p/2 \ge 1$. Leeson suggests the quoted $-100.8$\,dBc/Hz is closed-loop: check it.
  \item A two-stage PPF plus CML buffers gives I/Q within about 2$^\circ$; every degree is 1.11\,ps and every 1\% of duty cycle is 4\,ps at the half-rate samplers.
\end{itemize}
\end{takeaways}
